% Selected simulation evidence regenerated by the unified branch.

\subsubsection{Current terminal distribution and latent-state diagnostics}

The current Team~8 diagnostic simulation uses the 2 September calibration,
4,096 paths, 260 steps and the GLD level 402.78 USD/share. Its constant
annual rate is 4.0770\%, as recorded in the diagnostic table. This standalone
state-law illustration differs from the corporate adapter, which restarts at
Barrick's realised-price anchor and uses integral-preserving NSS forward rates
and 8,192 paths. Neither the levels nor the path percentiles are interchangeable.

\begin{table}[H]
\centering
\scriptsize
\setlength{\tabcolsep}{2.45pt}
\begin{tabular}{@{}lrrrrrrrr@{}}
\toprule
Model & Mean & Std & P05 & P50 & P95 & Return skew & Ex. kurt. & Mean jumps \\
\midrule
GBM / BS & 490.549 & 289.852 & 160.071 & 423.839 & 1036.482 & 1.743 & 5.708 & 0.000 \\
Heston & 495.256 & 308.314 & 164.068 & 421.222 & 1087.352 & 2.133 & 9.293 & 0.000 \\
Bates & 496.912 & 310.039 & 161.053 & 421.499 & 1098.107 & 2.069 & 7.779 & 1.259 \\
Full Bates--Hawkes & 494.830 & 300.874 & 164.799 & 425.021 & 1063.213 & 1.901 & 6.477 & 5.206 \\
\bottomrule
\end{tabular}
\caption{Terminal GLD price and simple-return diagnostics regenerated from the
current Team 8 calibration: 4,096 five-year paths, 260 steps and a fixed seed. The jump
models mainly alter tail shape and event counts; these levels are not Barrick
values and are not inputs to the later 8,192-path corporate comparison.}
\label{tab:t8-terminal-path-stats}
\end{table}

\begin{figure}[p]
    \centering
    \includegraphics[width=0.94\textwidth]{img/team8_empirical/terminal_return_percentiles.png}
    \caption{All integer percentiles 0--100 of five-year simple GLD returns in
    the current Team 8 simulation. The symmetric-log scale retains the
    central distribution while exposing extreme endpoints; the full integer
    percentile tables remain in the source article rather than being repeated
    in the thesis.}
    \label{fig:t8-terminal-return-percentiles}
\end{figure}

\begin{figure}[p]
    \centering
    \includegraphics[width=0.98\textwidth]{img/team8_empirical/volatility_state_paths.png}
    \caption{Representative annualized stochastic-volatility paths under
    Heston, Bates and Full Bates--Hawkes on a common scale. The plot verifies
    that all three richer engines carry a simulated variance state; the jump
    mechanism does not replace stochastic variance.}
    \label{fig:t8-volatility-state-paths}
\end{figure}

\begin{figure}[p]
    \centering
    \includegraphics[width=0.68\textwidth]{img/team8_empirical/bates_poisson_jump_paths.png}

    \vspace{3mm}
    \includegraphics[width=0.68\textwidth]{img/team8_empirical/hawkes_jump_paths.png}
    \caption{Representative jump-state diagnostics. Top: Bates retains a
    constant Poisson intensity while cumulative counts evolve independently.
    Bottom: Full Bates--Hawkes intensity rises after arrivals and decays toward
    its baseline, generating clustered event counts. These path displays
    explain the state difference behind the tail comparison; they are not
    additional forecast observations.}
    \label{fig:t8-jump-state-comparison}
\end{figure}
